\documentclass[11pt,a4paper]{article}
\usepackage[utf8]{inputenc}
\usepackage{amsmath,amssymb,amsthm}
\usepackage{graphicx}
\usepackage{booktabs}
\usepackage{hyperref}
\usepackage[margin=2.5cm]{geometry}
\usepackage{xcolor}

\hypersetup{
  colorlinks=true,
  linkcolor=blue,
  citecolor=blue,
  urlcolor=blue,
  pdfauthor={David Alfyorov},
  pdftitle={DYN-2: Dynamical Normalization of the CJ Bridge Formula}
}

\newtheorem{theorem}{Theorem}
\newtheorem{proposition}{Proposition}
\newtheorem{remark}{Remark}

\title{DYN-2: Dynamical Normalization of the CJ Bridge Formula\\[6pt]
\large Numerical investigation and time-reversal mechanism}
\author{David Alfyorov}
\date{April 2026}

\begin{document}
\maketitle
\tableofcontents

\vspace{1em}
\noindent\textit{Credits: Igor Shnyukov (physical interpretation, analytical framework).}

% ============================================================
\section{Problem Statement}

The CJ bridge formula for a Poisson-sprinkled $d{=}4$ Alexandrov diamond reads
\begin{equation}\label{eq:bridge}
  \mathbb{E}[\mathrm{CJ}_N] = C_0\, N^{8/9}\, \mathcal{E}^2\, T^4 \,(1 + o(1)),
\end{equation}
where $C_0 = 32\pi^2/(3 \cdot 9! \cdot 45) \approx 6.447 \times 10^{-6}$
and $\mathcal{E}^2 = E_{ij}E^{ij}$ is the electric part of the Weyl tensor.

The coefficient $C_0 = 4\,C_\mathrm{AN}$, where $C_\mathrm{AN} = 8\pi^2/(3 \cdot 9! \cdot 45)$
is assembled from the analytical coefficient skeleton (Benincasa--Dowker normalization,
beta-overlap identities, angular-volume arithmetic).
The factor~4 decomposes as
\begin{equation}\label{eq:factor4}
  4 = \underbrace{2}_{\text{algebraic}} \times \underbrace{2}_{\text{dynamical}}.
\end{equation}

The \textbf{algebraic~2} is exact: since $\delta Y \approx g_- + g_+$ in the bulk,
$\delta Y^2 = 2\,g_s^2$ where $g_s = (g_- + g_+)/\sqrt{2}$ is the symmetric channel.
Hence $\mathrm{CJ} = 2\,M_{ss}$.

The \textbf{dynamical~2} is the claim
\begin{equation}\label{eq:dyn2}
  \frac{M_{ss}}{C_\mathrm{AN}\, N^{8/9}\, \mathcal{E}^2\, T^4} \;\longrightarrow\; 2
  \qquad\text{as } N \to \infty.
\end{equation}

This note reports the first high-$N$ investigation of the dynamical factor,
including spatial diagnostics that identify the time-reversal mechanism behind it.

% ============================================================
\section{Channel Decomposition}

Define the past and future curvature responses
\[
  g_- = \log_2(p_{\downarrow}^{\mathrm{curv}}) - \log_2(p_{\downarrow}^{\mathrm{flat}}),\qquad
  g_+ = \log_2(p_{\uparrow}^{\mathrm{curv}}) - \log_2(p_{\uparrow}^{\mathrm{flat}}),
\]
and the stratified covariances
\[
  M_{--} = Q_N[g_-^2],\quad M_{++} = Q_N[g_+^2],\quad M_{-+} = Q_N[g_-\,g_+],
\]
where $Q_N[F] = \sum_B w_B\,\mathrm{Cov}_B(|X|, F)$.

The symmetric and antisymmetric channel covariances are
\begin{align}
  M_{ss} &= \tfrac{1}{2}(M_{--} + 2\,M_{-+} + M_{++}), \label{eq:Mss}\\
  M_{aa} &= \tfrac{1}{2}(M_{--} - 2\,M_{-+} + M_{++}), \label{eq:Maa}\\
  M_{sa} &= \tfrac{1}{2}(M_{--} - M_{++}). \label{eq:Msa}
\end{align}

Defining the single-leg average $L = (M_{--}^{\mathrm{norm}} + M_{++}^{\mathrm{norm}})/2$
and the effective correlation
\begin{equation}\label{eq:rho}
  \rho_{\mathrm{eff}} = \frac{M_{-+}}{\sqrt{M_{--}\,M_{++}}},
\end{equation}
the dynamical normalization becomes
\begin{equation}\label{eq:constraint}
  M_{ss}^{\mathrm{norm}} = L \cdot (1 + \rho_{\mathrm{eff}}).
\end{equation}
The target $M_{ss}^{\mathrm{norm}} = 2$ is satisfied whenever $L(1 + \rho_{\mathrm{eff}}) = 2$.

% ============================================================
\section{Numerical Results}

\subsection{Phase~1: High-$N$ channel decomposition}

We measure the channel covariances using common random numbers (CRN)
on pp-wave backgrounds ($\varepsilon = 3$, $T = 1$) at $N = 3000$ to $20{,}000$,
with $M = 20$ trials per $N$.

\begin{table}[h]
\centering
\caption{Channel normalization data.  All quantities normalized by $C_\mathrm{AN}\,N^{8/9}\,\mathcal{E}^2\,T^4$.}
\label{tab:channels}
\begin{tabular}{rrrrr}
\toprule
$N$ & $M_{ss}^{\mathrm{norm}}$ & $\rho_{\mathrm{eff}}$ & $\mathrm{corr}(g_-,g_+)$ & $M_{aa}^{\mathrm{norm}}$ \\
\midrule
3\,000  & $1.79 \pm 0.10$ & $0.577 \pm 0.032$ & 0.499 & 0.534 \\
5\,000  & $1.99 \pm 0.10$ & $0.578 \pm 0.019$ & 0.556 & 0.550 \\
10\,000 & $1.93 \pm 0.05$ & $0.598 \pm 0.012$ & 0.597 & 0.504 \\
15\,000 & $2.08 \pm 0.06$ & $0.611 \pm 0.010$ & 0.628 & 0.507 \\
20\,000 & $1.87 \pm 0.04$ & $0.598 \pm 0.010$ & 0.635 & 0.479 \\
\bottomrule
\end{tabular}
\end{table}

\textbf{Key observations:}
\begin{enumerate}
\item $M_{ss}^{\mathrm{norm}}$ clusters around~2 at all tested~$N$ (mean $1.97 \pm 0.05$).
\item $\rho_{\mathrm{eff}}$ is positive and growing: $0.58 \to 0.60 \to 0.61$.
\item The pointwise Pearson correlation $\mathrm{corr}(g_-,g_+)$ grows more steadily: $0.50 \to 0.56 \to 0.60 \to 0.63 \to 0.64$.
\item $M_{aa}^{\mathrm{norm}}$ is persistently nonzero ($\approx 0.50$), definitively excluding the equalization hypothesis $M_{--} = M_{-+} = M_{++}$.
\end{enumerate}

\subsection{Phase~2: Spatial diagnostics}

At $N = 10{,}000$ with $M = 40$ trials, we decompose the pointwise correlation $\rho(g_-,g_+)$
into 20 bins along the normalized depth coordinate~$\zeta$.

\begin{table}[h]
\centering
\caption{Phase~2 decomposition at $N = 2000, 5000, 10000$.}
\label{tab:phase2}
\begin{tabular}{rrrrrr}
\toprule
$N$ & $L$ & $\rho_{\mathrm{eff}}$ & $L(1{+}\rho)$ & $\cos(g_-,g_+)$ & median$(g_+/g_-)$ \\
\midrule
2\,000  & 1.259 & 0.559 & 1.99 & 0.461 & 0.75 \\
5\,000  & 1.272 & 0.568 & 2.00 & 0.559 & 0.85 \\
10\,000 & 1.239 & 0.566 & 1.94 & 0.600 & 0.89 \\
\bottomrule
\end{tabular}
\end{table}

\textbf{Spatial profile} (Fig.~\ref{fig:profile}):
$\rho(\zeta)$ is bell-shaped, peaking at $\zeta \approx 0.45{-}0.55$ (diamond center)
with $\rho_{\mathrm{peak}} \approx 0.63$ at $N = 10{,}000$, and dropping to $\rho \lesssim 0.2$
near the boundaries.  The asymmetry $\langle|g_- - g_+|\rangle/\langle|g_-| + |g_+|\rangle$
shows the complementary pattern: minimal ($\approx 0.45$) at the center, maximal ($\approx 0.80$)
at the boundaries.

% ============================================================
\section{Time-Reversal Mechanism}

\subsection{Physical picture}

The Minkowski diamond $\mathcal{D}(T)$ has an approximate time-reversal symmetry at its center:
the reflection $t \mapsto -t$ maps past light cones to future light cones and vice versa.
Under a small metric perturbation (here a pp-wave), the curvature response at a point~$x$
depends on the geometry of its past cone ($g_-$) and future cone ($g_+$).

At the diamond center, the past and future cones are images of each other under time reversal.
Both ``see'' the same spacetime curvature, so both respond similarly: $g_-(x) \approx g_+(x)$.
Near the boundaries, one cone is truncated by the diamond boundary while the other extends
into the bulk, breaking the symmetry.

The spatial profile confirms this picture:
\begin{itemize}
\item Center ($\zeta \approx 0.5$): $\rho \approx 0.63$, asymmetry $\approx 0.45$
\item Boundary ($\zeta \lesssim 0.1$ or $\zeta \gtrsim 0.9$): $\rho \lesssim 0.2$, asymmetry $\approx 0.8$
\end{itemize}

\subsection{Consequence for the factor~2}

If $g_- = g_+$ exactly at all bulk points, then $M_{-+} = M_{--} = M_{++}$ (so $\rho_{\mathrm{eff}} = 1$)
and
\[
  M_{ss} = \tfrac{1}{2}(M_{--} + 2M_{--} + M_{--}) = 2\,M_{--},
\]
giving the dynamical~2 as the ratio $M_{ss}/M_{--} = 2$.

In practice, $g_- \neq g_+$ even at the center (the pp-wave breaks time-reversal),
but the positive correlation $\rho_{\mathrm{eff}} \approx 0.6$ at $N = 20{,}000$
contributes a factor $(1 + 0.6) = 1.6$, while the single-leg normalization
$L \approx 1.25$ provides the remaining factor, so $L(1{+}\rho) \approx 2.0$.

Whether $\rho_{\mathrm{eff}} \to 1$ in the strict continuum limit
($N \to \infty$) remains open.  The data is consistent with slow convergence
(the bulk central value grows from $\approx 0.52$ at $N{=}3000$ to $\approx 0.66$ at $N{=}20{,}000$),
but the rate is too slow to confirm extrapolation from the tested range alone.

% ============================================================
\section{Analytical Framework: Even/Odd Decomposition}

\subsection{Channel decomposition}

Define the symmetric and antisymmetric responses:
\[
  s(x) = \frac{g_-(x) + g_+(x)}{\sqrt{2}}\,,\qquad
  a(x) = \frac{g_-(x) - g_+(x)}{\sqrt{2}}\,.
\]
Then $g_-(x)\,g_+(x) = \frac{1}{2}(s^2 - a^2)$, and the stratified
covariances satisfy
\begin{equation}\label{eq:Mpm-decomp}
  M_{-+} = \frac{M_{ss} - M_{aa}}{2}\,.
\end{equation}
Positive cross-correlation ($M_{-+} > 0$) is therefore equivalent to
$M_{ss} > M_{aa}$: the symmetric channel dominates the antisymmetric one.

\subsection{Kernel-response model}

For a pp-wave perturbation $h_{ij} = \varepsilon^2 f(u)\,\delta_{ij}$,
the curvature responses admit the continuum representation
\[
  g_\pm(x) \approx \varepsilon^2
  \int_0^\infty K_N(x,s)\,f(u_x \pm s)\,ds + B_\pm^\partial(x)\,,
\]
where $K_N$ is a common kernel for past/future in the bulk and
$B_\pm^\partial$ are boundary corrections.  Taylor-expanding $f(u_x \pm s)$:
\[
  g_\pm(x) = c(x) \pm b(x)\,,
\]
where $c(x)$ collects even moments of $K_N$ (common response) and
$b(x)$ collects odd moments (directional asymmetry from the wave
propagation and boundary truncation).  Then
\[
  g_-(x)\,g_+(x) = c(x)^2 - b(x)^2\,.
\]
In the bulk center, $c^2 \gg b^2$ (time-reversal approximate symmetry),
giving strong positive correlation.  Near boundaries, $B_a^\partial$ grows
(one cone truncated more than the other), reducing $\rho$.

This explains the bell-shaped spatial profile of $\rho(\zeta)$:
the even channel dominates at the center and the odd channel grows
at the boundaries.

\subsection{Exact formula for $\rho_{\mathrm{eff}}$}

\begin{proposition}[Exact rho formula]\label{prop:rho-exact}
Define $r_N := M_{aa}^{\mathrm{norm}} / M_{ss}^{\mathrm{norm}}$.
When $M_{sa}/M_{ss}$ is small,
\begin{equation}\label{eq:rho-exact}
  \rho_{\mathrm{eff}} \approx \frac{1 - r_N}{1 + r_N}\,.
\end{equation}
Therefore $\rho_{\mathrm{eff}} \to 1$ if and only if $r_N \to 0$,
i.e., the antisymmetric channel vanishes relative to the symmetric one.
\end{proposition}

\begin{proof}
From~\eqref{eq:Mpm-decomp} and $M_{--} = (M_{ss} + M_{aa} + 2M_{sa})/2$,
$M_{++} = (M_{ss} + M_{aa} - 2M_{sa})/2$:
\[
  \rho_{\mathrm{eff}}
  = \frac{M_{ss} - M_{aa}}
  {\sqrt{(M_{ss}+M_{aa})^2 - 4M_{sa}^2}}\,.
\]
For $|M_{sa}| \ll M_{ss}$, this reduces to $(1-r_N)/(1+r_N)$.
\end{proof}

Numerical check from our data:
\begin{center}
\begin{tabular}{rrrrr}
\toprule
$N$ & $r_N = M_{aa}/M_{ss}$ & $\rho_{\text{predicted}}$ & $\rho_{\text{measured}}$ & Deviation \\
\midrule
3\,000  & 0.298 & 0.541 & 0.577 & 4\% \\
5\,000  & 0.276 & 0.567 & 0.578 & 2\% \\
10\,000 & 0.261 & 0.586 & 0.598 & 2\% \\
15\,000 & 0.244 & 0.608 & 0.611 & 1\% \\
20\,000 & 0.256 & 0.592 & 0.598 & 1\% \\
\bottomrule
\end{tabular}
\end{center}

\subsection{Assessment: $\rho_{\mathrm{eff}} \to 1$ or $R < 1$?}

Five arguments support the hypothesis $\rho_{\mathrm{eff}} \to R < 1$
with $R \approx 0.6$:

\begin{enumerate}
\item \textbf{Global time-reversal breaking.}
  The pp-wave depends on the retarded coordinate $u = t - z$.
  Under $t \mapsto -t$: $u \mapsto -v$, so the wave does not map to itself.
  Exact time-reversal symmetry is absent even in the continuum.

\item \textbf{Global observables.}
  $p_\downarrow(x)$ counts paths from the bottom vertex; $p_\uparrow(x)$
  from the top.  For $t \neq 0$, the past and future distances to the
  vertices differ, introducing inherent asymmetry.

\item \textbf{Fixed bulk mask.}
  The boundary cut $\zeta = 0.15$ is fixed in continuum units and does
  not shrink with~$N$.  If the odd contribution persists in the fixed bulk,
  it need not vanish as $N \to \infty$.

\item \textbf{Stable ratio.}
  The data show $L \approx 1.25$, $\rho \approx 0.58$, with product
  $L(1{+}\rho) \approx 2.0$ stable across all~$N$.  This is consistent
  with a finite limit $R \approx 0.6$, $L \to 2/(1{+}R) \approx 1.25$.

\item \textbf{Central asymmetry not negligible.}
  Even at $\zeta = 0.5$, the pointwise asymmetry is $\approx 0.45$ (45\%),
  not $\approx 0$.  This does not suggest the antisymmetric channel is
  nearly extinct.
\end{enumerate}

\begin{remark}[Open question: $|X|$-weight localization]
  One mechanism that could still drive $\rho \to 1$ despite nonvanishing
  pointwise asymmetry: if the CJ weight $|X|$ preferentially samples
  the central region (where $\rho$ is highest) with increasing~$N$,
  the weighted $\rho_{\mathrm{eff}}$ could approach~1 even as the
  unweighted correlation saturates.  Whether such localization occurs
  is an open question requiring analysis of the $|X|$-distribution
  as a function of~$N$.
\end{remark}

% ============================================================
\section{Established Results}

\begin{enumerate}
\item $M_{ss}^{\mathrm{norm}} = 2 \pm 5\%$ at all tested $N \in [3000, 20000]$.

\item The past-future curvature responses $g_-$ and $g_+$ are positively correlated,
with $\rho_{\mathrm{eff}}$ growing from $0.58$ to $0.61$ over $N = 3000$ to $20{,}000$.

\item The correlation is spatially structured: bell-shaped, peaking at the diamond center,
consistent with time-reversal symmetry as the mechanism.

\item The equalization hypothesis ($M_{--} = M_{-+} = M_{++}$) is definitively excluded:
$M_{aa}^{\mathrm{norm}} \approx 0.50$ is persistently nonzero.

\item The factor~2 arises from the product $L \cdot (1 + \rho_{\mathrm{eff}}) = 2$,
where both $L$ (single-leg normalization) and $\rho_{\mathrm{eff}}$ (past-future correlation)
evolve with~$N$ in a way that keeps their product fixed.
\end{enumerate}

% ============================================================
\section{Open Questions}

\begin{enumerate}
\item \textbf{Does $\rho_{\mathrm{eff}} \to 1$ or $R < 1$?}
  Current data ($N \leq 20{,}000$) supports $R \approx 0.6$ more
  than $\rho \to 1$ (see Section~5.4).  However, $|X|$-weight
  localization could still drive $\rho \to 1$ (Remark in Section~5.4).
  Resolving this requires either $N \gg 20{,}000$ data or a proof
  that $M_{aa}/M_{ss} \to 0$.

\item \textbf{Analytical derivation of $M_{ss}^{\mathrm{norm}} = 2$.}
  The even/odd decomposition (Section~5) provides a framework but not
  a proof.  A complete derivation requires the continuum limit of the
  resolvent $(I - A)^{-1}$ acting on path counts, which involves the
  nonlocal structure of the Hasse transitive reduction.

\item \textbf{Dimension dependence.}
  The CJ bridge formula involves $N^{2d/(2d+1)}$; the even/odd channel
  decomposition might have a different ratio $r = M_{aa}/M_{ss}$ in $d \neq 4$.

\item \textbf{$|X|$-weight localization.}
  Does the CJ weight $|X| = |Y_{\mathrm{flat}} - \langle Y_{\mathrm{flat}}\rangle|$
  concentrate on the central region with increasing~$N$?  If so, the
  weighted $\rho_{\mathrm{eff}}$ could approach~1 even with stable
  pointwise asymmetry.
\end{enumerate}

% ============================================================
\section{Verification Summary}

\begin{itemize}
\item Phase~1: 6 values of~$N$ ($3000$--$20{,}000$), 20 CRN trials each, $\varepsilon = 3$. Total: 120 MC trials.
\item Phase~2: 3 values of~$N$ ($2000$, $5000$, $10{,}000$), 40 CRN trials each. Total: 120 MC trials. 20-bin spatial profiles.
\item Total computation time: $\approx 2$ hours on RTX~3090~Ti.
\item Data archived in \texttt{dynamical2\_highN.json} and \texttt{dynamical2\_rho\_diagnostic.json}.
\item 4 publication figures in \texttt{analysis/figures/dyn2/}.
\end{itemize}

\end{document}
